\newpage

\phantomsection  % 自动生成无名锚点
%
\renewcommand{\bibname}{\protect\hyperlink{:}{Reference}}
\addtocontents{toc}{\protect\hypertarget{:}{}} %用\hypertarget实现反向跳转
\addcontentsline{toc}{chapter}{Reference}
\thispagestyle{fancy} % 强制当前页（第一页）使用 fancy 样式
\bibliographystyle{ieeetr}
\normalem	% 去掉下划线
\nocite{*}	% 引用未被引用的文献
%
\bibliography{refence/refences}  % 引用参考文献的文件

%添加空白页，根据需要灵活取舍
\clearpage
\null
\clearpage

%相关研究\cite{ref1}和\cite{ref2,ref3}表明XXX。
%\begin{thebibliography}{99}  %99表示参考文献的最大个数为99
%	\thispagestyle{fancy} % 强制当前页（第一页）使用 fancy 样式、
%	
%	\bibitem{ref1}Zheng L, Wang S, Tian L, et al., Query-adaptive late fusion for image search and person re-identification, Proceedings of the IEEE Conference on Computer Vision and Pattern Recognition, 2015: 1741-1750.  
%	\bibitem{ref2}Arandjelović R, Zisserman A, Three things everyone should know to improve object retrieval, Computer Vision and Pattern Recognition (CVPR), 2012 IEEE Conference on, IEEE, 2012: 2911-2918.  
%	\bibitem{ref3}Lowe D G. Distinctive image features from scale-invariant keypoints, International journal of computer vision, 2004, 60(2): 91-110.  
%\end{thebibliography}